% paper3_4_supplemental.tex
% Supplemental Material for Papers 3 and 4:
%   Paper 3: "Dark Matter as Information Loss: Running G from the Gravitational Channel"
%   Paper 4: "Five Papers, One Equation: Sigma = D(rho_spacetime || rho_matter)"
% Author: Sheng-Kai Huang
% Compile with: pdflatex paper3_4_supplemental.tex

\documentclass[aps,prd,preprint,superscriptaddress,longbibliography]{revtex4-2}

\usepackage{amsmath}
\usepackage{amssymb}
\usepackage{amsfonts}
\usepackage{amsthm}
\usepackage{bm}
\usepackage{hyperref}
\usepackage{booktabs}
\usepackage{graphicx}

\newtheorem{theorem}{Theorem}
\newtheorem{proposition}[theorem]{Proposition}
\newtheorem{conjecture}[theorem]{Conjecture}
\newtheorem{definition}[theorem]{Definition}
\newtheorem{lemma}[theorem]{Lemma}

% Convenience macros (same as main texts)
\newcommand{\kB}{k_{\mathrm{B}}}
\newcommand{\Tr}{\mathrm{Tr}}
\newcommand{\id}{\mathrm{id}}
\newcommand{\Petz}{\mathcal{R}}
\newcommand{\chan}{\mathcal{N}}
\newcommand{\hilb}{\mathcal{H}}
\newcommand{\code}{\mathcal{C}}
\newcommand{\KL}{D}
\newcommand{\ket}[1]{|#1\rangle}
\newcommand{\bra}[1]{\langle#1|}
\newcommand{\braket}[2]{\langle#1|#2\rangle}
\newcommand{\abs}[1]{\left|#1\right|}
\newcommand{\norm}[1]{\left\|#1\right\|}
\newcommand{\tauasym}{\tau}
\newcommand{\GN}{G_{\mathrm{N}}}
\newcommand{\rs}{r_{\mathrm{s}}}
\newcommand{\LdS}{L_{\mathrm{dS}}}
\newcommand{\adS}{a_{\mathrm{dS}}}
\newcommand{\azero}{a_0}

\begin{document}

\title{Supplemental Material for Papers~3 and~4:\\
Dark Matter as Information Loss and the Unified $\Sigma$ Framework}

\author{Sheng-Kai Huang}
\email{akai@fawstudio.com}
\affiliation{Independent Researcher}

\date{\today}

\maketitle

\tableofcontents

%=======================================================================
%  PART I: PAPER 3 SUPPLEMENT
%=======================================================================

\part*{Part I: Paper 3 --- Dark Matter from the Gravitational Channel}

%=======================================================================
\section{S1. Running $G$ Derivation from the Gravitational Channel}
\label{sec:running_G}
%=======================================================================

\subsection{The scale-dependent gravitational channel}

The gravitational channel at momentum scale $k$ is defined as
\begin{equation}
\chan_{\mathrm{grav}}(k): \rho \mapsto \Tr_{E(k)}
\big[U(G(k))\,\rho\, U(G(k))^{\dagger}\big],
\label{eq:grav_channel}
\end{equation}
where $U(G(k))$ is the unitary evolution with the Hamiltonian
$H(G(k))$ and the environment trace represents gravitational
decoherence at scale~$k$.  The entropy production at scale~$k$ is
\begin{equation}
\Sigma(k) = \KL(\rho\|\sigma) - \KL\!\big(\chan_{\mathrm{grav}}(k)(\rho)
\,\big\|\,\chan_{\mathrm{grav}}(k)(\sigma)\big).
\label{eq:Sigma_k}
\end{equation}

By the data processing inequality (DPI), $\Sigma(k)\ge 0$ for all~$k$.
For scales $k_{\mathrm{IR}} < k_{\mathrm{UV}}$, coarse-graining from
UV to IR constitutes an additional channel, giving the monotonicity
\begin{equation}
\Sigma(k_{\mathrm{IR}}) \ge \Sigma(k_{\mathrm{UV}}) \ge 0.
\label{eq:DPI_monotone}
\end{equation}
This is the Casini--Huerta result~\cite{CasiniTeste2017}: the
renormalization group (RG) flow is a quantum channel, and the QRE
decreases monotonically under the RG flow.

\subsection{Kumar's QFT first-principles running}

Following Kumar~\cite{Kumar2025}, the gravitational coupling runs
under the RG with anomalous dimension $\eta$:
\begin{equation}
\frac{d\ln G}{d\ln k} = -\eta,
\qquad
G(k) = \GN\left(\frac{k_*}{k}\right)^{\!\eta}
\quad\text{for }k \ll k_*,
\label{eq:running_G}
\end{equation}
where $k_*$ is the crossover momentum and $\GN = G(k_*)$ is the
measured Newton constant.

The modified Newtonian potential in momentum space is
\begin{equation}
\tilde{\Phi}(k) = -\frac{4\pi G(k)\,M}{k^2}
= -\frac{4\pi\GN M\, k_*^{\eta}}{k^{2+\eta}}.
\label{eq:Phi_k}
\end{equation}

Fourier-transforming to position space in $d=3$ spatial dimensions:
\begin{equation}
\Phi(r) = -M \int \frac{d^3 k}{(2\pi)^3}\,
\frac{4\pi G(k)}{k^2}\,e^{i\bm{k}\cdot\bm{r}}.
\label{eq:Phi_FT}
\end{equation}

\subsection{Why $\eta=1$ is the marginal case}

The integral~\eqref{eq:Phi_FT} with $G(k)\propto k^{-\eta}$ yields
three qualitatively distinct behaviors:

\begin{center}
\begin{tabular}{lll}
\toprule
$\eta$ & $\Phi(r)$ behavior & Physical status \\
\midrule
$\eta < 1$ & $r^{\eta-1}$ (decaying) & IR-irrelevant \\
$\eta = 1$ & $\ln r$ (logarithmic) & \textbf{Marginal, scale-invariant} \\
$\eta > 1$ & $r^{\eta-1}$ (growing) & Violates solar-system constraints \\
\bottomrule
\end{tabular}
\end{center}

For $\eta=1$, the running becomes $G(k) = \GN(1 + k_*/k)$, and the
explicit position-space results are:
\begin{align}
\Phi(r) &= -\frac{\GN M}{r}
+ \frac{2\GN M k_*}{\pi}\,\ln\!\left(\frac{r}{r_0}\right),
\label{eq:Phi_eta1}\\
F(r) &= -\frac{\GN M}{r^2} - \frac{2\GN M k_*}{\pi\, r},
\label{eq:F_eta1}\\
v^2(r) &= r|F(r)| = \frac{\GN M}{r}
+ \frac{2\GN M k_*}{\pi}.
\label{eq:v2}
\end{align}

At $r \gg r_c = 1/k_*$, the second term dominates and
\begin{equation}
v^2(r) \approx V_0^2 = \frac{2\GN M_{\mathrm{bar}}\, k_*}{\pi}
= \mathrm{const},
\label{eq:flat_v}
\end{equation}
yielding a \emph{flat rotation curve}.

\subsection{Connection to $\Sigma_{\mathrm{grav}}$}

Using $\Sigma_{\mathrm{grav}} = 2|\Phi|/c^2$ (Paper~2) with the
running potential~\eqref{eq:Phi_eta1}:
\begin{equation}
\Sigma_{\mathrm{grav}}(r) = \frac{\rs}{r}
+ \frac{4k_*\rs}{\pi}\ln\!\left(\frac{r}{r_0}\right),
\label{eq:Sigma_galactic}
\end{equation}
where $\rs = 2\GN M/c^2$.  This is \emph{not} an ad hoc addition.
It is the same $\Sigma = -\ln(-g_{00})$ from Paper~2, evaluated with
the quantum-corrected potential that follows from the RG-improved
gravitational channel.

\subsection{SPARC galaxy rotation curve data}

Three SPARC galaxies fitted by Kumar~\cite{Kumar2025} with $\eta=1$
give a \emph{universal} crossover scale:

\begin{center}
\begin{tabular}{lccc}
\toprule
Galaxy & $M_{\mathrm{bar}}\,[M_\odot]$ & $V_0\,[\mathrm{km/s}]$
  & $k_*\,[\mathrm{kpc}^{-1}]$ \\
\midrule
S:700624 & $1.01\times 10^{12}$ & $271.4\pm 2.6$
  & $(2.66\pm 0.05)\times 10^{-2}$ \\
S:700916 & $5.85\times 10^{11}$ & $210.4\pm 2.4$
  & $(2.76\pm 0.06)\times 10^{-2}$ \\
S:705253 & $3.55\times 10^{11}$ & $158.9\pm 1.9$
  & $(2.60\pm 0.06)\times 10^{-2}$ \\
\bottomrule
\end{tabular}
\end{center}

The values of $k_*$ are consistent across a factor $\sim 3$ in
baryonic mass, giving $r_c \approx 37\,\mathrm{kpc}$.  NFW dark
matter halos have no analogous universal scale.

Gubitosi, Piattella, and Casarini~\cite{Gubitosi2024} performed
MCMC fits to 100 SPARC galaxies using the RGGR model
($G(\mu) = G_0/[1 + \nu\ln(\mu^2/\mu_0^2)]$) and found that the
running-$G$ framework is \emph{competitive with NFW}: 47~galaxies
prefer RGGR, 40~prefer NFW, and 13~are inconclusive.

%=======================================================================
\section{S2. Observer-Dependent $\Sigma_{\mathrm{DM}}$ Calculation}
\label{sec:observer_DM}
%=======================================================================

\subsection{Classical and quantum observer algebras}

\begin{definition}[Observer algebras]
For a gravitating system with total algebra $\mathcal{A}_{\mathrm{total}}$,
define:
\begin{itemize}
\item $\mathcal{A}_{\mathrm{classical}}$: the subalgebra of
  electromagnetic observables (photon frequencies, luminosities,
  positions from EM detection).
\item $\mathcal{A}_{\mathrm{quantum}} = \mathcal{A}_{\mathrm{total}}$:
  the full algebra including gravitational entanglement structure.
\end{itemize}
\end{definition}

By construction, $\mathcal{A}_{\mathrm{classical}} \subset
\mathcal{A}_{\mathrm{quantum}}$.  The restriction map
$\chan_{\mathrm{cl}} = E_{\mathcal{A}_{\mathrm{classical}}}$ is a
quantum channel (conditional expectation).

\subsection{Computing the two entropy productions}

By the DPI/monotonicity theorem (Paper~1, Theorem~1 of the
observer-dependent framework):
\begin{equation}
\Sigma_{\mathrm{classical}} \ge \Sigma_{\mathrm{quantum}}.
\label{eq:monotonicity}
\end{equation}

For a point mass $M$ in the weak-field regime:
\begin{align}
\Sigma_{\mathrm{quantum}}(r) &= \frac{\rs}{r},
\label{eq:Sigma_quantum}\\
\Sigma_{\mathrm{classical}}(r) &= \frac{\rs}{r}
+ \frac{4k_*\rs}{\pi}\ln\!\left(\frac{r}{r_0}\right)
+ \frac{r}{\LdS},
\label{eq:Sigma_classical}
\end{align}
where the additional terms arise because the classical observer
cannot access the gravitational entanglement that would reduce the
apparent information loss.

\subsection{Dark matter as the observer gap}

\begin{definition}[Dark matter entropy production]
\begin{equation}
\Sigma_{\mathrm{DM}}(r) \equiv
\Sigma_{\mathrm{classical}}(r) - \Sigma_{\mathrm{quantum}}(r)
= \frac{4k_*\rs}{\pi}\ln\!\left(\frac{r}{r_0}\right)
+ \frac{r}{\LdS}.
\label{eq:Sigma_DM}
\end{equation}
\end{definition}

The corresponding temporal asymmetry parameter for dark matter is
\begin{equation}
\tauasym_{\mathrm{DM}} = 1 - \exp(-\Sigma_{\mathrm{DM}}/2).
\label{eq:tau_DM}
\end{equation}

\subsection{Numerical estimate at the solar radius}

For the Milky Way ($M \approx 6\times 10^{10}\,M_\odot$) at
$r = 8\,\mathrm{kpc}$:
\begin{align}
\Sigma_{\mathrm{Newton}} &= \frac{\rs}{r}
\approx 3\times 10^{-6}, \nonumber\\
\Sigma_{\mathrm{running}} &\approx \frac{4k_*\rs}{\pi}
\ln\!\left(\frac{8}{37}\right) \approx -1\times 10^{-6}, \nonumber\\
\Sigma_{\mathrm{DM}} &\approx 3.5\times 10^{-7}, \nonumber\\
\tauasym_{\mathrm{DM}} &\approx 1.7\times 10^{-7}.
\label{eq:tau_solar}
\end{align}

This tiny but nonzero $\tauasym_{\mathrm{DM}}$ is the origin of the
galactic ``dark matter'' at this radius: it represents the information
inaccessible to a classical (electromagnetic) observer.

%=======================================================================
\section{S3. Crooks $\to$ $\azero$ Derivation: Eight Routes}
\label{sec:crooks_a0}
%=======================================================================

We systematically attempt to derive the MOND acceleration scale
$\azero \approx 1.2\times 10^{-10}\,\mathrm{m/s}^2$ from the
Crooks fluctuation theorem~\cite{Crooks1999} within the $\tauasym$
framework.

\subsection{Numerical benchmarks}

With $H_0 = 70\,\mathrm{km\,s^{-1}\,Mpc^{-1}}
= 2.27\times 10^{-18}\,\mathrm{s}^{-1}$:
\begin{equation}
\begin{aligned}
cH_0 &= 6.80\times 10^{-10}\,\mathrm{m/s}^2,\\
cH_0/(2\pi) &= 1.08\times 10^{-10}\,\mathrm{m/s}^2,\\
cH_0/6 &= 1.13\times 10^{-10}\,\mathrm{m/s}^2,\\
\azero\,\text{(MOND)} &= 1.20\pm 0.02\times 10^{-10}\,\mathrm{m/s}^2.
\end{aligned}
\label{eq:benchmarks}
\end{equation}

\subsection{Route 1: Naive $\Sigma=1$ transition (FAILS)}

Setting $\Sigma_{\mathrm{grav}} = \rs/r = 1$ gives
$a = c^4/(4GM)$, which is mass-dependent.  For the Milky Way,
$a \sim 10^2\,\mathrm{m/s}^2$, a factor $10^{12}$ above $\azero$.
The naive transition identifies the event horizon, not the MOND
scale.

\subsection{Route 2: $\Sigma_{\mathrm{grav}} = \Sigma_{\mathrm{cosmo}}$
  (FAILS)}

Balancing $\rs/r = r^2/r_H^2$ gives
$r_{\mathrm{tr}} = (\rs\,r_H^2)^{1/3}$ and
$a_{\mathrm{tr}} = (GM)^{1/3}H_0^{4/3}/2^{2/3}$.  This is still
mass-dependent and $\sim 10^{-14}\,\mathrm{m/s}^2$.

\subsection{Route 3: Verlinde in $\tauasym$ language (WORKS, imported)}

Verlinde's emergent gravity~\cite{Verlinde2016} gives
\begin{equation}
g_D(r) = \sqrt{\frac{cH_0\,g_B(r)}{6}},
\label{eq:Verlinde}
\end{equation}
with transition $g_D = g_B$ at $\azero = cH_0/6 \approx
1.13\times 10^{-10}\,\mathrm{m/s}^2$.

In $\tauasym$ language: in the Newtonian regime ($g\gg\azero$),
$\Sigma \approx \rs/r$ and $\tauasym \approx \rs/(2r)$; in the MOND
regime ($g\ll\azero$), $\Sigma_{\mathrm{eff}}$ grows as $\ln r$
from the displaced de Sitter entanglement, and $\tauasym$ exceeds
the Newtonian prediction.  However, the factor~6 is imported from
Verlinde's specific geometric assumptions ($V/A = r/3$, elastic
factor~$1/2$) rather than derived from Crooks.

\subsection{Route 4: Unruh--de Sitter matching (FAILS by $\sim 6\times$)}

Setting $T_U = T_{\mathrm{dS}}$ gives $a = cH_0 \approx
6.8\times 10^{-10}\,\mathrm{m/s}^2$, a factor $\sim 6$ too large.

\subsection{Route 5: Entropic force matching (circular)}

Setting the Verlinde dark force equal to the Newtonian force
reproduces $\azero = cH_0/6$, but the calculation imports Verlinde's
specific formula, making it circular.

\subsection{Route 6: Volume-law entropy with Crooks (INCOMPLETE)}

The displaced Crooks entropy
$\Sigma_{\mathrm{displaced}} \sim S_V\,\epsilon^2/2$ involves
$S_V\propto Mc^2/(\hbar H_0)$ and $\epsilon\sim\rs/r$.  The product
is mass-dependent.  To extract a universal $\azero$, one needs
Verlinde's prescription for counting degrees of freedom on holographic
screens.

\subsection{Route 7: Geometric factor derivation (circular)}

The factor $6 = 3\times 2$ in Verlinde's $\azero = cH_0/6$ comes
from: $3$ from $V/A = r/3$ (spherical geometry), and $2$ from the
quadratic QRE near equilibrium $\KL(\rho+\delta\|\rho) \approx
\tfrac{1}{2}\Tr(\delta^2/\rho)$.  These factors are
\emph{geometrically motivated} from the $\tauasym$ framework but not
derived from Crooks.

\subsection{Route 8: KMS--Crooks connection (BEST)}

The de Sitter vacuum satisfies the KMS condition with thermal
period $\beta_{\mathrm{dS}} = 2\pi/H_0$.  The spatial thermal
wavelength is $\lambda_{\mathrm{dS}} = c\beta_{\mathrm{dS}}
= 2\pi c/H_0$, and the corresponding acceleration is
\begin{equation}
\azero = \frac{c^2}{\lambda_{\mathrm{dS}}}
= \frac{c^2}{2\pi c/H_0} = \frac{cH_0}{2\pi}
\approx 1.08\times 10^{-10}\,\mathrm{m/s}^2,
\label{eq:a0_KMS}
\end{equation}
within $10\%$ of the empirical value.

The physical logic:
\begin{enumerate}
\item The KMS condition is exact for the de Sitter equilibrium state.
\item The Crooks theorem requires a thermal state; KMS is the
  mathematical definition of such a state in QFT.
\item The modular flow defines ``thermal time,'' and Crooks entropy
  production $\Sigma$ is measured in these units.
\item The MOND transition $\Sigma\sim 1$ occurs when the
  gravitational acceleration matches the de Sitter thermal scale.
\end{enumerate}

This reframes MOND as \emph{marginal irreversibility}: the transition
from $\Sigma \gg 1$ (Newtonian, strongly irreversible) to $\Sigma\ll 1$
(deep MOND, nearly reversible).

\subsection{Summary of routes}

\begin{center}
\begin{tabular}{clcc}
\toprule
Route & Method & Result & Verdict \\
\midrule
1 & Naive $\Sigma=1$ & Mass-dep. & FAILS \\
2 & $\Sigma_{\mathrm{grav}}=\Sigma_{\mathrm{cosmo}}$ & Mass-dep. & FAILS \\
3 & Verlinde in $\tauasym$ & $cH_0/6$ & Works (imported) \\
4 & Unruh--dS match & $cH_0$ & FAILS ($\times 6$) \\
5 & Entropic force & $cH_0/6$ & Circular \\
6 & Volume-law Crooks & Incomplete & --- \\
7 & Geometric factor & $cH_0/6$ & Circular \\
8 & KMS--Crooks & $cH_0/(2\pi)$ & \textbf{Best} \\
\bottomrule
\end{tabular}
\end{center}

All routes agree: $\azero \sim cH_0/\mathcal{O}(1)$, with the
$\mathcal{O}(1)$ factor between $1$ and $2\pi$.

%=======================================================================
\section{S4. Bullet Cluster: Entanglement Rigidity}
\label{sec:bullet}
%=======================================================================

\subsection{The challenge}

The Bullet Cluster (1E~0657-558) shows an offset between the
gravitational lensing center (tracing total mass) and the X-ray
center (tracing hot gas).  In $\Lambda$CDM, this arises because
collisionless dark matter passes through while gas is ram-pressure
stripped.  Any modified gravity theory must explain this offset.

\subsection{Timescale estimates}

In the $\tauasym$ framework, the ``dark component'' is
gravitational entanglement of the stellar/galactic degrees of
freedom.  Define two timescales:
\begin{align}
t_{\mathrm{ent}} &= \frac{r_{\mathrm{gal}}}{c}
\approx \frac{30\,\mathrm{kpc}}{c}
\approx 10^5\,\mathrm{yr},
\label{eq:t_ent}\\
t_{\mathrm{coll}} &\approx \frac{D_{\mathrm{sep}}}{v_{\mathrm{rel}}}
\approx \frac{700\,\mathrm{kpc}}{4700\,\mathrm{km/s}}
\approx 1.5\times 10^8\,\mathrm{yr},
\label{eq:t_coll}
\end{align}
where $t_{\mathrm{ent}}$ is the light-crossing time for the galaxy's
entanglement region and $t_{\mathrm{coll}}$ is the cluster crossing
time.

Since $t_{\mathrm{coll}} \gg t_{\mathrm{ent}}$, the entanglement
structure can respond (approximately) adiabatically to the collision
dynamics.  The entanglement entropy is anchored to the stellar
distribution (which is collisionless), not to the gas (which is
collisional).

\subsection{Why entanglement stays with galaxies}

The running-$G$ contribution to $\Sigma$ depends on the matter
distribution through $M_{\mathrm{bar}}(r)$.  Stars and dark
matter (in the conventional picture) are collisionless; gas is
collisional.  In the $\tauasym$ framework:
\begin{equation}
\Sigma_{\mathrm{total}}(r) = \underbrace{\frac{\rs^{\mathrm{stars}}}{r}
+ \frac{4k_*\rs^{\mathrm{stars}}}{\pi}\ln r}_{\text{follows stars}}
+ \underbrace{\frac{\rs^{\mathrm{gas}}}{r}}_{\text{follows gas}}.
\label{eq:bullet}
\end{equation}

The logarithmic (running-$G$) correction is dominated by the stellar
mass distribution, which passes through the collision unimpeded.
The gas contribution to the Newtonian term is subdominant
($M_{\mathrm{gas}} \ll M_{\mathrm{stars}}$ in the subcluster
galaxies).

\subsection{Quantitative lensing-gas offset bound}

The predicted lensing center is at the center of mass of
$\Sigma_{\mathrm{total}}$.  The gas displacement $\Delta x_{\mathrm{gas}}
\approx 200$--$300\,\mathrm{kpc}$ (observed) produces a shift in
the lensing center of order
\begin{equation}
\Delta x_{\mathrm{lens}} \approx \frac{M_{\mathrm{gas}}}{M_{\mathrm{total}}}
\,\Delta x_{\mathrm{gas}} \sim 0.1\,\Delta x_{\mathrm{gas}}
\sim 20\text{--}30\,\mathrm{kpc},
\label{eq:offset}
\end{equation}
consistent with the observed lensing-gas offset of $\sim 150\,\mathrm{kpc}$
only if the running-$G$ contribution provides the remaining
``collisionless'' component.  This requires
$\Sigma_{\mathrm{running}} \gtrsim 4\,\Sigma_{\mathrm{Newton}}$ at
the cluster scale, which remains to be quantitatively verified.

%=======================================================================
\section{S5. Crooks Interpolation Function}
\label{sec:interpolation}
%=======================================================================

\subsection{Derivation from $P(\Sigma)/P(-\Sigma) = e^{\Sigma}$}

The Crooks relation gives $P_F(\Sigma)/P_R(-\Sigma) = e^{\Sigma}$.
If we identify $\Sigma = a/\azero$ (the local acceleration in units
of the MOND scale), then the ratio of forward to reverse
probabilities is $\exp(a/\azero)$, and the ``efficiency'' of the
gravitational process is
\begin{equation}
\mu_{\mathrm{Crooks}}(x) = 1 - e^{-x},
\qquad x = \frac{a}{\azero} = \frac{g_{\mathrm{bar}}}{\azero}.
\label{eq:mu_Crooks}
\end{equation}

Properties:
\begin{itemize}
\item $x \gg 1$: $\mu \approx 1$ (Newtonian regime).
\item $x \ll 1$: $\mu \approx x$ (deep MOND, linear).
\item The transition is at $x\sim 1$ ($a\sim\azero$).
\end{itemize}

\subsection{Comparison with standard MOND interpolation functions}

\begin{center}
\begin{tabular}{lcccc}
\toprule
Function & $\mu(x)$ & $x=0.1$ & $x=1$ & $x=3$ \\
\midrule
Standard & $x/(1+x)$ & 0.091 & 0.500 & 0.750 \\
Simple & $x/\sqrt{1+x^2}$ & 0.100 & 0.707 & 0.949 \\
RAR (empirical) & $1-e^{-\sqrt{x}}$ & 0.272 & 0.632 & 0.826 \\
Crooks-inspired & $1-e^{-x}$ & 0.095 & 0.632 & 0.950 \\
\bottomrule
\end{tabular}
\end{center}

The Crooks-inspired form transitions faster than the ``standard''
function but matches the ``simple'' function well for $x \ge 1$.
Its behavior at $x = 1$ ($\mu = 1-1/e \approx 0.632$) coincides
with the empirical RAR interpolation $1-e^{-\sqrt{x}}$ evaluated at
$x=1$.

\subsection{Rotation curve predictions}

The effective gravitational acceleration is
\begin{equation}
g_{\mathrm{obs}} = \frac{g_{\mathrm{bar}}}{\mu(g_{\mathrm{bar}}/\azero)}
= \frac{g_{\mathrm{bar}}}{1-\exp(-g_{\mathrm{bar}}/\azero)}.
\label{eq:gobs}
\end{equation}

For a galaxy with baryonic mass profile $M_{\mathrm{bar}}(r)$, the
circular velocity is
$v^2(r) = r\,g_{\mathrm{obs}}(r)$, reproducing flat rotation curves
in the regime $g_{\mathrm{bar}} \ll \azero$.  Direct comparison with
SPARC data requires fitting $\Upsilon_*$ (stellar mass-to-light
ratio) as the sole free parameter.

%=======================================================================
\section{S6. Dark Energy from Cosmic Information (CosMIn)}
\label{sec:dark_energy}
%=======================================================================

\subsection{Padmanabhan's argument}

Padmanabhan~\cite{Padmanabhan2017} argues that the cosmological
constant is determined by the requirement that the total cosmic
information be finite:
\begin{equation}
I_{\mathrm{CosMIn}} = \int_0^{\infty}
(N_{\mathrm{sur}}(t) - N_{\mathrm{bulk}}(t))\,dt < \infty,
\label{eq:cosmin}
\end{equation}
where $N_{\mathrm{sur}} = A/(l_P^2)$ and $N_{\mathrm{bulk}}$ are
the surface and bulk degrees of freedom of the Hubble sphere.

\subsection{Translation to the $\tauasym$ framework}

In our language:
\begin{equation}
I_{\mathrm{CosMIn}} = \int_0^{\infty} \tauasym(t)\,dt < \infty,
\label{eq:cosmin_tau}
\end{equation}
where $\tauasym(t)$ is the cosmological temporal asymmetry at time~$t$.
The requirement that this integral converges demands
$\tauasym(t) \to 0$ sufficiently fast as $t\to\infty$, which requires
accelerated expansion ($\Lambda > 0$).

\subsection{Connection to the cosmological horizon}

At the de Sitter future ($\Lambda$-dominated era), the Hubble sphere
approaches $r_H = c/H_{\Lambda}$ with
$H_{\Lambda} = \sqrt{\Lambda/3}$.  The horizon entropy is
$S_H = \pi r_H^2 / (\GN\hbar)$, and the Kodama-vector $\tauasym$
profile is
\begin{equation}
\tauasym(r) = 1 - \sqrt{1 - H_\Lambda^2 r^2/c^2},
\quad r < c/H_\Lambda.
\label{eq:tau_dS}
\end{equation}

This goes from $\tauasym = 0$ at the observer to $\tauasym = 1$ at
the horizon, confirming that the cosmological horizon is the
maximum-$\tauasym$ boundary.

%=======================================================================
%  PART II: PAPER 4 SUPPLEMENT
%=======================================================================

\part*{Part II: Paper 4 --- Grand Unification via
$\Sigma = \KL(\rho_{\mathrm{st}}\|\rho_{\mathrm{m}})$}

%=======================================================================
\section{S7. Background QRE on FRW}
\label{sec:FRW_background}
%=======================================================================

\subsection{Cai--Kim apparent horizon thermodynamics}

For flat FRW with scale factor $a(t)$ and Hubble parameter
$H = \dot{a}/a$, the apparent horizon is at $r_A = 1/H$.  Its
thermodynamic quantities are~\cite{CaiKim2005}:
\begin{align}
T_A &= \frac{H}{2\pi},
\quad
S_A = \frac{\pi}{\GN H^2},
\quad
E = \frac{1}{2\GN H}.
\label{eq:CK_thermo}
\end{align}

\subsection{$\KL_{\mathrm{background}} = 0$ implies the Friedmann equations}

Define the spacetime thermal state and matter state by
\begin{align}
\rho_{\mathrm{st}} &= \frac{1}{Z_A}
\exp\!\Big(-\frac{2\pi}{H}\,M(r_A)\Big),\\
\rho_{\mathrm{m}} &= \text{state encoding }T_{ab}\text{ within }r_A.
\end{align}

The QRE between them:
$\KL(\rho_{\mathrm{st}}\|\rho_{\mathrm{m}})
= \beta_A\,\Delta E - \Delta S$.  At equilibrium
(Friedmann equations satisfied), $\KL = 0$: the spacetime and matter
states are in ``entanglement equilibrium''~\cite{Jacobson2016}.

\subsection{Detailed derivation}

Applying the Clausius relation $-dE = T_A\,dS_A$ at the apparent
horizon:

\emph{Energy}: $E = \rho\cdot(4\pi/3)r_A^3 = 4\pi\rho/(3H^3)$.

\emph{Differentiating}: Using $\dot{\rho} = -3H(\rho+p)$:
\begin{equation}
\dot{E} = \frac{4\pi}{3}\left[\frac{\dot{\rho}}{H^3}
- \frac{3\rho\dot{H}}{H^4}\right]
= \frac{4\pi}{3}\left[\frac{-3H(\rho+p)}{H^3}
- \frac{3\rho\dot{H}}{H^4}\right].
\end{equation}

\emph{Entropy}: $S_A = \pi/(\GN H^2)$, so
$\dot{S}_A = -2\pi\dot{H}/(\GN H^3)$.

\emph{Clausius relation}: $-\dot{E} = T_A\,\dot{S}_A$ gives
\begin{equation}
4\pi\!\left[\frac{\rho+p}{H^2} + \frac{\rho\dot{H}}{H^3}\right]
= \frac{\dot{H}}{\GN H^2}.
\end{equation}

This yields:
\begin{align}
\dot{H} &= -4\pi\GN(\rho+p)
\quad\text{(second Friedmann equation)},
\label{eq:Friedmann2}\\
H^2 &= \frac{8\pi\GN}{3}\rho
\quad\text{(first Friedmann equation, from the constraint)}.
\label{eq:Friedmann1}
\end{align}

\subsection{Consistency check: $\Sigma_{\mathrm{grav}} = 0$ on FRW}

In comoving coordinates, $g_{00} = -1$, so
$\Sigma_{\mathrm{grav}} = -\ln(-g_{00}) = 0$.  This is consistent:
a comoving observer experiences zero gravitational entropy production.
The cosmological arrow of time arises entirely from expansion
(increasing scale factor), not from gravitational information loss.

%=======================================================================
\section{S8. First-Order QRE Perturbation}
\label{sec:first_order}
%=======================================================================

\subsection{The Aalsma--Bak modular Hamiltonian perturbation}

Following Aalsma and Bak~\cite{AalsmaBak2025}, the modular
Hamiltonian for the FRW causal diamond, perturbed by a Newtonian
potential $\Phi$, becomes
\begin{equation}
\delta K = 2\,S_A\,\Phi,
\label{eq:deltaK}
\end{equation}
where $S_A = \pi/(\GN H^2)$ is the de Sitter (apparent horizon) entropy.

\subsection{Entanglement first law: $\delta\KL^{(1)}=0$}

The first-order variation of the QRE is given by the entanglement
first law~\cite{Faulkner2014,BlancoCasini2013}:
\begin{equation}
\delta\KL^{(1)} = \delta\langle K_A\rangle - \delta S_A.
\label{eq:ent_first_law}
\end{equation}

The modular Hamiltonian perturbation gives
$\delta\langle K_A\rangle = 2S_A\,\Phi$.

The entropy perturbation is
$\delta S_A = -S_A\,\delta_m$, where $\delta_m = \delta\rho/\rho$
is the matter density contrast (from the shift in the apparent
horizon location).

Setting $\delta\KL^{(1)} = 0$ (entanglement equilibrium at first
order):
\begin{equation}
2\Phi + \delta_m = 0
\quad\Longrightarrow\quad
\Phi = -\frac{1}{2}\,\delta_m = -\frac{\delta\rho}{2\rho}.
\label{eq:Poisson}
\end{equation}

This is the \textbf{Poisson equation in the sub-Hubble limit}:
$k^2\Phi = -4\pi\GN a^2\rho\,\delta_m$, which for modes well inside
the horizon ($k\gg H$) reduces to $\Phi \approx -\delta_m/2$.

\subsection{Gauge choices}

The above is derived in Newtonian gauge:
$ds^2 = -(1+2\Phi)\,dt^2 + a^2(1-2\Psi)\,\delta_{ij}\,dx^i\,dx^j$.
For GR without anisotropic stress, $\Phi = \Psi$.  The Bardeen
potentials are gauge-invariant and coincide with $\Phi,\Psi$ in
Newtonian gauge.  In synchronous gauge, the calculation requires
expressing $\delta K$ and $\delta S_A$ in terms of the synchronous
potentials $h$ and $\eta$.

\subsection{What appears and what does not}

At first order in the standard QRE formalism:
\begin{itemize}
\item \textbf{Present}: scalar metric perturbations ($\Phi,\Psi$)
  and matter perturbations ($\delta\rho,\delta p,\delta q$).
\item \textbf{Absent}: any independent ``observer field'' or
  Khronon-like scalar.  The Kodama vector, which plays the role of
  the observer, is determined by the metric---not an independent
  degree of freedom.
\end{itemize}

%=======================================================================
\section{S9. Modular Khronon Emergence}
\label{sec:khronon}
%=======================================================================

\subsection{The modular time perturbation}

The modular Hamiltonian $K$ generates a flow with modular parameter~$s$.
The ``modular time'' associated with metric perturbation $\Phi$ is
\begin{equation}
\delta T_{\mathrm{mod}} = \frac{\delta K}{2\pi H^2/(2\pi)}
= \frac{2S_A\,\Phi}{H} = -\frac{2\Phi}{H}.
\label{eq:delta_T_mod}
\end{equation}

This perturbation defines a ``displaced foliation''---the
constant-modular-time surfaces are shifted from the constant-cosmic-time
surfaces by $\delta T_{\mathrm{mod}}$.

\subsection{Dispersion relation}

On a Minkowski background, the modular time perturbation
$\delta T_{\mathrm{mod}}$ satisfies the constraint
$\delta\KL^{(1)} = 0$, which ties $\delta T_{\mathrm{mod}}$ to
$\Phi$.  The perturbation does \emph{not} satisfy a wave equation;
instead, it is determined algebraically by the Poisson
equation~\eqref{eq:Poisson}.  Its effective dispersion relation is
\begin{equation}
\omega = 0 \quad\text{(non-propagating)},
\label{eq:cs0}
\end{equation}
because the modular-time perturbation is a \emph{constraint mode},
not a propagating degree of freedom.

In a matter-dominated universe, the Newtonian potential is constant
($\dot{\Phi}=0$), so $\delta T_{\mathrm{mod}}$ is frozen---it grows
in place without oscillating, exactly like pressureless dust (CDM).

\subsection{Why $c_s \neq 0$ in the radiation era}

During radiation domination, the gravitational potential decays
($\dot{\Phi}\neq 0$), so $\delta T_{\mathrm{mod}}$ is driven.
The effective sound speed
$c_s^2 = \delta p/\delta\rho$ for the modular perturbation is
\begin{equation}
c_s^2 \approx \begin{cases}
0 & \text{matter era (}w=0\text{)},\\
\mathcal{O}(k/\mu)^2 & \text{radiation era, subleading}.
\end{cases}
\label{eq:cs_rad}
\end{equation}

For all CMB-relevant scales ($k\ll\mu$ if $\mu^{-1}\sim\mathrm{kpc}$),
$c_s^2\approx 0$ to excellent approximation.

\subsection{Effective mass}

The Khronon mass parameter
$\mu$ sets the scale at which the non-propagating behavior breaks
down.  From the Blanchet--Skordis theory~\cite{BlanchetSkordis2024}:
\begin{equation}
\mu_{\mathrm{eff}} \sim 2\pi H_0
\approx 1.4\times 10^{-17}\,\mathrm{s}^{-1}.
\label{eq:mu_eff}
\end{equation}

In position space, $\mu^{-1}\sim 223\,\mathrm{kpc}$, consistent with
the crossover between MOND and cosmological regimes.

%=======================================================================
\section{S10. Observer Entanglement Equilibrium (OEE)}
\label{sec:OEE}
%=======================================================================

\subsection{Formal statement}

\begin{definition}[Observer Entanglement Equilibrium]
The physical state $\omega$ satisfies OEE if the QRE is stationary
with respect to variations in the metric \emph{and} the observer's
4-velocity:
\begin{equation}
\frac{\delta\KL}{\delta g_{ab}} = 0,
\qquad
\frac{\delta\KL}{\delta u^a} = 0,
\label{eq:OEE}
\end{equation}
subject to $g_{ab}u^a u^b = -1$ (enforced by a Lagrange multiplier
$\lambda$).
\end{definition}

The first equation gives the Einstein equations (Dorau-Much/Jacobson).
The second is new: it determines the observer's 4-velocity as a
\emph{dynamical} field.

\subsection{Derivation of the Khronon field equation from OEE}

The QRE depends on $u^a$ through the modular Hamiltonian
$K[u] = (2\pi/\kappa_u)\int T_{ab}\,u^a\,d\Sigma^b$.
Varying with respect to $u^a$ (with the constraint
$u^a u_a = -1$) yields
\begin{equation}
(g_{ab}+u_a u_b)\,\nabla^b\!\left(\frac{\partial K}{\partial u^a}\right)
= 0,
\label{eq:khronon_eom}
\end{equation}
which, for $u^a = -(c/Q)\nabla^a T$ (the hypersurface-orthogonal
form), reduces to the Khronon field equation with kinetic function
\begin{equation}
K(Q) = \mu^2(Q-1)^2.
\label{eq:KQ}
\end{equation}

\subsection{Connection to Petz optimality}

The Petz optimality argument: the kinetic function $K(Q)$ must have a
quadratic minimum at the background value $Q_0 = 1$, because:
\begin{enumerate}
\item Petz optimality requires the channel to be maximally recoverable.
\item Maximal recoverability implies minimal entropy production for perturbations.
\item Minimal entropy production requires perturbations not to propagate
  ($c_s = 0$).
\item $c_s = 0$ requires $K(Q)$ to have a quadratic minimum at the
  background: $K(Q) \sim \mu^2(Q-Q_0)^2$.
\end{enumerate}

This selects exactly the Blanchet--Skordis kinetic
term~\cite{BlanchetSkordis2024}.

\subsection{Why OEE gives $K(Q) = \mu^2(Q-1)^2$}

The variation $\delta\KL/\delta u^a = 0$ at the background
($u^a$ = comoving) gives $K'(Q_0) = 0$.  Combined with stability
($K''(Q_0) > 0$), the leading-order kinetic function is
\begin{equation}
K(Q) = K(Q_0) + \tfrac{1}{2}K''(Q_0)(Q-Q_0)^2 + \cdots
= \mu^2(Q-1)^2 + \mathcal{O}((Q-1)^3).
\end{equation}

%=======================================================================
\section{S11. $\tauasym\to$ AeST Mapping}
\label{sec:AeST}
%=======================================================================

\subsection{Complete AeST action}

The Skordis--Zlosnik AeST action~\cite{SkordisZlosnik2021} is
\begin{equation}
S = \int d^4x\,\frac{\sqrt{-g}}{16\pi\tilde{G}}
\Big[R - \frac{K_B}{2}F_{\mu\nu}F^{\mu\nu}
+ 2(2-K_B)J^\mu\nabla_\mu\phi
- (2-K_B)Y - \mathcal{F}(Y,Q)
- \lambda(A^\mu A_\mu + 1)\Big]
+ S_{\mathrm{m}},
\label{eq:AeST}
\end{equation}
where $F_{\mu\nu} = \nabla_\mu A_\nu - \nabla_\nu A_\mu$,
$J^\mu = A^\nu\nabla_\nu A^\mu$ (aether acceleration),
$Y = D^\mu\phi\,D_\mu\phi$ (spatial scalar kinetic term with
$D_\mu\phi = (\delta^\nu_\mu + A^\nu A_\mu)\nabla_\nu\phi$),
and $Q = A^\mu\nabla_\mu\phi$ (temporal scalar derivative).

\subsection{The identification: $u^a = A^\mu$}

The structural identification is exact:

\begin{center}
\begin{tabular}{lll}
\toprule
$\tauasym$ framework & AeST / Khronon & Math.\ object \\
\midrule
Observer $u^a$ & Aether $A^\mu$ & Unit timelike vector \\
Time function $T$ & Khronon $\tau$ & Scalar foliation \\
$u^a u_a = -1$ & $A^\mu A_\mu = -1$ & Same constraint \\
$\Sigma = -\ln(-g_{00})$ & $\phi_s$ & Gravitational scalar \\
DPI monotonicity & MOND interpolation & Running of channel \\
\bottomrule
\end{tabular}
\end{center}

\subsection{What the $\tauasym$ framework determines vs.\ what is free}

\textbf{Determined by $\tauasym$}:
\begin{enumerate}
\item The existence of a dynamical unit timelike field $u^a$ (from
  observer-dependence of $\tauasym$).
\item The quadratic-minimum form of the kinetic function $K(Q)$
  (from Petz optimality; Section~S10).
\item The sign constraints $c_a = c_1+c_4 \ge 0$ (from DPI/energy
  conditions).
\item The constraint $c_1+c_3 = 0$ (from $c_T = c$;
  GW170817~\cite{GW170817}).
\end{enumerate}

\textbf{Free parameters (not determined)}:
\begin{enumerate}
\item The coupling constants $c_1$ and $c_4$ (two free parameters).
\item The mass scale $\mu$ (sets the DM energy density).
\item The MOND acceleration scale $\azero$ (entered via $\mathcal{F}(Y,Q)$).
\item The DBI completion (behavior away from the quadratic minimum).
\end{enumerate}

\subsection{The stealth black hole tension}

AeST admits ``stealth'' black hole solutions with exact
Schwarzschild geometry~\cite{SkordisVokrouhlicky2024}, while the
$\tauasym$ framework predicts the exponential metric
$g_{00} = -e^{-\rs/r}$ (Paper~2).  These are distinct and
observationally distinguishable.  This tension is real and must be
acknowledged; its resolution requires determining whether the Petz
saturation ansatz or the AeST field equations take precedence in the
strong-field regime.

%=======================================================================
\section{S12. Five Limits of the Master Equation}
\label{sec:five_limits}
%=======================================================================

The master equation
$\Sigma = \KL(\rho_{\mathrm{st}}\|\rho_{\mathrm{m}})$ reduces to
known physics under different boundary conditions.

\subsection{Paper 1 limit: Quantum information}

\emph{Setting}: Flat spacetime ($g_{\mu\nu} = \eta_{\mu\nu}$), CPTP
map $\chan$, finite-dimensional Hilbert space.

$\Sigma_{\mathrm{grav}} = 0$ (no gravitational contribution).
The remaining $\Sigma$ is
\begin{equation}
\Sigma_{\mathrm{QI}} = \KL(\rho\|\sigma) - \KL(\chan(\rho)\|\chan(\sigma)).
\label{eq:limit1}
\end{equation}

The Petz recovery bound $F\ge\exp(-\Sigma_{\mathrm{QI}}/2)$, the
equivalence chain $\tauasym=0\Leftrightarrow\Sigma=0
\Leftrightarrow$ quantum Markov $\Leftrightarrow$ Knill--Laflamme,
and the retrodiction Landauer principle all follow.

\emph{Precision}: EXACT.  No approximations.

\subsection{Paper 2 limit: Strong-field gravity}

\emph{Setting}: Static, spherically symmetric, asymptotically flat,
$r\sim\rs$.

The gravitational redshift channel gives
$\Sigma_{\mathrm{grav}} = -\ln(-g_{00})$.  For the exponential metric
(exact Petz saturation):
$g_{00} = -\exp(-\rs/r)$, hence $\Sigma = \rs/r$.  Three
independent first-principles routes confirm this
(modular flow, gravitational Landauer, pure-loss bosonic channel).

\emph{Precision}: Steps 1--3 DERIVED; exact saturation CONJECTURED.

\subsection{Paper 3 limit: Galactic dynamics}

\emph{Setting}: Weak field, $r\sim\mathrm{kpc}$, RG-improved channel.

The DPI applied to the RG flow (Casini--Huerta~\cite{CasiniTeste2017})
gives monotonicity $\Sigma(k_{\mathrm{IR}})\ge\Sigma(k_{\mathrm{UV}})$.
With $\eta=1$ running (Kumar~\cite{Kumar2025}):
\begin{equation}
\Sigma_{\mathrm{weak}}(r) = \frac{\rs}{r}
+ \frac{4k_*\rs}{\pi}\ln\!\left(\frac{r}{r_0}\right).
\label{eq:limit3}
\end{equation}

This produces flat rotation curves, the radial acceleration
relation, and the baryonic Tully-Fisher relation.

\emph{Precision}: Running $G$ KNOWN (asymptotic safety);
$\eta=1$ PHENOMENOLOGICAL (not derived from $\tauasym$ framework).

\subsection{Paper 5 limit: Observer-dependent time}

\emph{Setting}: General system, observer $O$ with algebra
$\mathcal{A}_O\subset\mathcal{A}_{\mathrm{total}}$.

The observer's effective channel is
$\chan_O:\mathcal{A}_{\mathrm{total}}\to\mathcal{A}_O$ (conditional
expectation).  Observer-dependent entropy production:
\begin{equation}
\Sigma_O = \KL(\omega\|\omega_{\mathrm{ref}})
- \KL(\chan_O(\omega)\|\chan_O(\omega_{\mathrm{ref}})),
\end{equation}
with monotonicity: $O_1\le O_2\Rightarrow\Sigma_{O_1}\ge\Sigma_{O_2}$.

\emph{Precision}: Monotonicity PROVED (DPI).

\subsection{Friedmann limit: Cosmological expansion}

\emph{Setting}: Homogeneous, isotropic FRW, $r\sim c/H_0$.

At the background level, $\KL = 0$ yields the Friedmann equations
(Section~S7).  At first order, $\delta\KL^{(1)}=0$ yields the
Poisson equation (Section~S8).  The Kodama-vector $\tauasym$ profile
$\tauasym(r) = 1 - \sqrt{1-H^2r^2/c^2}$ interpolates from $\tauasym=0$
at the observer to $\tauasym=1$ at the apparent horizon.

\emph{Precision}: Background PROVED (Cai-Kim~\cite{CaiKim2005});
perturbations NEW SYNTHESIS (Aalsma--Bak~\cite{AalsmaBak2025}).

%=======================================================================
\section{S13. Why $\Omega_{\mathrm{DM}}$ Cannot Be Derived}
\label{sec:Omega_DM}
%=======================================================================

\subsection{The fundamental obstruction}

The JRSWW bound $F\ge\exp(-\Sigma/2)$ is a \emph{universal}
information-theoretic statement.  It holds for every quantum channel,
every pair of states, and every Hilbert space dimension.  It knows
nothing about:
\begin{itemize}
\item the specific matter content of the universe,
\item the baryon-to-photon ratio,
\item Big Bang nucleosynthesis,
\item freeze-out temperatures.
\end{itemize}

Meanwhile, $\Omega_{\mathrm{DM}}h^2 \approx 0.12$ is a
\emph{contingent} fact about our universe.
\textbf{A universal bound cannot determine a contingent parameter.}

\subsection{Summary of 10 failed attempts}

\begin{center}
\begin{tabular}{clc}
\toprule
\# & Approach & Verdict \\
\midrule
1 & $\Sigma=-\ln(g_{00})$ at fixed $r/r_H$ & DEAD END (arbitrary $r$) \\
2 & $\tauasym = \Omega_{\mathrm{DM}}$ & DEAD END (tautological) \\
3 & Numerology & DEAD END \\
4 & Padmanabhan $N_{\mathrm{sur}}$ vs $N_{\mathrm{bulk}}$ & Predicts $\Lambda$, not $\Omega_{\mathrm{DM}}$ \\
5 & Cumulative $\Sigma$ at $r_H$ & CIRCULAR ($\Sigma\sim\Omega_m$) \\
6 & Bekenstein bound & DEAD END ($S/S_{\mathrm{holo}}\sim 10^{-18}$) \\
7 & $\Sigma=1$ at horizon & Suggestive ($\tauasym=0.39$ vs $\Omega_m=0.31$) \\
8 & Holographic DM (Trivedi-Scherrer) & Indirect connection only \\
9 & $N_{\mathrm{bulk}}/N_{\mathrm{sur}}$ ratio & CIRCULAR \\
10 & Baryogenesis entropy & DEAD END \\
\bottomrule
\end{tabular}
\end{center}

\subsection{The correct analogy}

$\Omega_{\mathrm{DM}}$ is to the $\tauasym$ framework as the mass
of the Earth is to $F=ma$.  Newton's second law provides the
\emph{framework} (how things respond to forces) but not the
\emph{initial conditions} (the mass of Earth).  Similarly,
$\Sigma = \KL(\rho_{\mathrm{st}}\|\rho_{\mathrm{m}})$ provides the
\emph{structure} of gravity and dark matter but not the
\emph{amplitude} $\Omega_{\mathrm{DM}}h^2 = 0.12$, which requires
cosmological initial conditions.

%=======================================================================
\section{S14. Comparison with Other Unification Programs}
\label{sec:comparison}
%=======================================================================

\subsection{String theory / M-theory}

\begin{center}
\begin{tabular}{lcc}
\toprule
Feature & String theory & $\tauasym$ framework \\
\midrule
Starting point & Strings/branes & QRE ($\Sigma$) \\
UV complete? & Yes & No \\
Background independent? & No (perturbative) & Yes (algebraic) \\
Predictions & Extra dims, SUSY & GW echoes, flat curves \\
Testable (2026)? & No & Yes \\
\bottomrule
\end{tabular}
\end{center}

The $\tauasym$ framework is an \emph{organizational principle}, not a
competitor to microscopic theories.  If string theory provides the
UV completion, $\tauasym$ may provide the macroscopic organizational
structure---analogous to thermodynamics vis-\`{a}-vis statistical
mechanics.

\subsection{Loop quantum gravity}

LQG quantizes geometry itself (spin networks, discrete area spectrum).
The $\tauasym$ framework does not quantize geometry but uses
information loss as the fundamental quantity.  A potential connection:
the discrete area spectrum in LQG would quantize
$\Sigma_{\mathrm{grav}}$ in Planck units, providing a UV cutoff.

\subsection{Entropic gravity (Verlinde)}

Verlinde~\cite{Verlinde2016} derives gravity from holographic
entropy (area + volume law).  The $\tauasym$ framework shares the
philosophy (gravity = entropy production) but uses QRE rather than
force-based entropy arguments.  Verlinde's volume-law contribution
may be the cosmological-scale limit of our $\Sigma$, while our
$\Sigma_{\mathrm{grav}} = \rs/r$ is the strong-field limit.
Ghari and Haghi~\cite{GhariHaghi2026} find Verlinde favored over
MOND at $5.2\sigma$ for dwarf spheroidals.

\subsection{ER = EPR}

The $\tauasym$ interpretation: a wormhole connecting two regions is
a system for which the total observer (accessing both ends) has
$\tauasym = 0$ (closed system, no arrow of time).  An observer
at one end only has $\tauasym > 0$ (the other end appears
thermalized).  This is the ER = EPR picture translated to $\tauasym$
language.  Kaplan, Mittal, Penington, and Susskind (2024) proved
this operationally for typical states.

\subsection{It from Qubit / Holographic QEC}

The Petz recovery map appears in both frameworks:
\begin{itemize}
\item In holographic QEC, the Petz map reconstructs bulk operators
  from boundary data (Cotler et al., PRX 2019).
\item In the $\tauasym$ framework, the Petz map is the optimal
  retrodiction (Paper~1).
\end{itemize}

These are the \emph{same} mathematical object in different contexts.
The $\tauasym$ framework extends beyond AdS/CFT (which requires a
conformal boundary) because QRE is well-defined in any algebraic QFT.

\subsection{Summary: what type of theory $\tauasym$ is}

The $\tauasym$ framework is not a microscopic theory (like string or
LQG).  It is a \emph{structural theory}---universal predictions that
hold for any microscopic completion satisfying the DPI.  Its value:
\begin{enumerate}
\item Explains \emph{why} Petz/QRE structures appear across domains.
\item Provides universal, model-independent predictions.
\item Connects QI, GR, cosmology, and philosophy of time.
\item Gives a computable definition of ``arrow of time'' at all scales.
\end{enumerate}

%=======================================================================
%  REFERENCES
%=======================================================================

\begin{thebibliography}{99}

\bibitem{CasiniTeste2017}
H.~Casini, E.~Teste, and G.~Torroba,
``Relative Entropy and the RG Flow,''
J. High Energy Phys.\ \textbf{2017}, 089 (2017);
arXiv:1611.00016.

\bibitem{Kumar2025}
N.~Kumar,
``Marginal IR Running of Gravity as a Natural Explanation for
Dark Matter,''
Phys.\ Lett.\ B \textbf{871}, 140008 (2025);
arXiv:2509.05246.

\bibitem{Gubitosi2024}
G.~Gubitosi, O.~F.~Piattella, and L.~Casarini,
``Phenomenology of Renormalization Group Improved Gravity from
SPARC Galaxies,''
Phys.\ Rev.\ D \textbf{110}, 124014 (2024);
arXiv:2403.00531.

\bibitem{Verlinde2016}
E.~P.~Verlinde,
``Emergent Gravity and the Dark Universe,''
SciPost Phys.\ \textbf{2}, 016 (2017);
arXiv:1611.02269.

\bibitem{GhariHaghi2026}
A.~Ghari and H.~Haghi,
(2026); arXiv:2601.01715.

\bibitem{Crooks1999}
G.~E.~Crooks,
``Entropy Production Fluctuation Theorem and the Nonequilibrium
Work Relation for Free Energy Differences,''
Phys.\ Rev.\ E \textbf{60}, 2721 (1999);
arXiv:cond-mat/9901352.

\bibitem{Padmanabhan2017}
T.~Padmanabhan and H.~Padmanabhan,
``Cosmic Information, the Cosmological Constant and the Amplitude
of Primordial Perturbations,''
Phys.\ Lett.\ B \textbf{773}, 81 (2017);
arXiv:1703.06144.

\bibitem{CaiKim2005}
R.-G.~Cai and S.~P.~Kim,
``First Law of Thermodynamics and Friedmann Equations of
Friedmann--Robertson--Walker Universe,''
J. High Energy Phys.\ \textbf{2005}, 050 (2005);
hep-th/0501055.

\bibitem{Jacobson2016}
T.~Jacobson,
``Entanglement Equilibrium and the Einstein Equation,''
Phys.\ Rev.\ Lett.\ \textbf{116}, 201101 (2016);
arXiv:1505.04753.

\bibitem{AalsmaBak2025}
W.~Aalsma and D.~Bak,
``Modular Hamiltonian and Cosmological Perturbations,''
Phys.\ Rev.\ D \textbf{112}, 026017 (2025);
arXiv:2503.04886.

\bibitem{Faulkner2014}
T.~Faulkner, A.~Lewkowycz, and J.~Maldacena,
``Quantum Corrections to Holographic Entanglement Entropy,''
J. High Energy Phys.\ \textbf{2013}, 074 (2013);
arXiv:1307.2892.

\bibitem{BlancoCasini2013}
D.~D.~Blanco, H.~Casini, L.-Y.~Hung, and R.~C.~Myers,
``Relative Entropy and Holography,''
J. High Energy Phys.\ \textbf{2013}, 015 (2013);
arXiv:1305.3182.

\bibitem{BlanchetSkordis2024}
L.~Blanchet and C.~Skordis,
``A New Relativistic Theory for Modified Newtonian Dynamics,''
J. Cosmol.\ Astropart.\ Phys.\ \textbf{2024}, 040 (2024);
arXiv:2404.06584.

\bibitem{SkordisZlosnik2021}
C.~Skordis and T.~Zlosnik,
``New Relativistic Theory for Modified Newtonian Dynamics,''
Phys.\ Rev.\ Lett.\ \textbf{127}, 161302 (2021);
arXiv:2007.00082.

\bibitem{SkordisVokrouhlicky2024}
C.~Skordis and D.~Vokrouhlick\'{y},
``Stealth Black Holes in AeST,''
(2024); arXiv:2412.15395.

\bibitem{GW170817}
B.~P.~Abbott \emph{et al.} (LIGO/Virgo),
``Tests of General Relativity with GW170817,''
Phys.\ Rev.\ Lett.\ \textbf{119}, 161101 (2017).

\bibitem{Bianconi2025}
G.~Bianconi,
``Gravity from Entropy,''
Phys.\ Rev.\ D \textbf{111}, 066001 (2025);
arXiv:2408.14391.

\bibitem{DorauMuch2025}
P.~Dorau and A.~Much,
``From Quantum Relative Entropy to the Semiclassical Einstein
Equations,''
Phys.\ Rev.\ Lett.\ \textbf{136}, 091602 (2025);
arXiv:2510.24491.

\bibitem{Milgrom1983}
M.~Milgrom,
``A Modification of the Newtonian Dynamics as a Possible
Alternative to the Hidden Mass Hypothesis,''
Astrophys.\ J.\ \textbf{270}, 365 (1983).

\bibitem{McGaugh2016}
S.~S.~McGaugh, F.~Lelli, and J.~M.~Schombert,
``Radial Acceleration Relation in Rotationally Supported Galaxies,''
Phys.\ Rev.\ Lett.\ \textbf{117}, 201101 (2016);
arXiv:1609.05917.

\bibitem{Lelli2016}
F.~Lelli, S.~S.~McGaugh, and J.~M.~Schombert,
``SPARC: Mass Models for 175 Disk Galaxies,''
Astron.\ J.\ \textbf{152}, 157 (2016);
arXiv:1606.09251.

\bibitem{BassoMazieroCeleri2025}
M.~L.~W.~Basso, J.~Maziero, and L.~C.~Celeri,
``Quantum Detailed Fluctuation Theorem in Curved Spacetimes,''
Phys.\ Rev.\ Lett.\ \textbf{134}, 050406 (2025);
arXiv:2405.03902.

\bibitem{Jacobson1995}
T.~Jacobson,
``Thermodynamics of Spacetime: The Einstein Equation of State,''
Phys.\ Rev.\ Lett.\ \textbf{75}, 1260 (1995);
gr-qc/9504004.

\bibitem{Petz1988}
D.~Petz,
``Sufficiency of Channels over von Neumann Algebras,''
Q.~J.~Math.\ \textbf{39}, 97 (1988).

\bibitem{Junge2018}
M.~Junge, R.~Renner, D.~Sutter, M.~M.~Wilde, and A.~Winter,
``Universal Recovery Maps and Approximate Sufficiency of Quantum
Relative Entropy,''
Ann.\ Henri Poincar\'{e} \textbf{19}, 2955 (2018);
arXiv:1509.07127.

\bibitem{ReuterWeyer2004}
M.~Reuter and H.~Weyer,
``Running Newton Constant, Improved Gravitational Actions, and
Galaxy Rotation Curves,''
Phys.\ Rev.\ D \textbf{70}, 124028 (2004);
hep-th/0410117.

\bibitem{BlanchetSkordis2025}
L.~Blanchet and C.~Skordis,
``The Khronon-Tensor Theory,''
(2025); arXiv:2507.00912.

\bibitem{Smolin2017}
L.~Smolin,
``MOND as a Regime of Quantum Gravity,''
Phys.\ Rev.\ D \textbf{96}, 083523 (2017);
arXiv:1704.00780.

\bibitem{Milgrom2020}
M.~Milgrom,
``The $a_0$-cosmology connection in MOND,''
(2020); arXiv:2001.09729.

\bibitem{Haubner2024}
M.~Haubner, F.~Lelli, \emph{et al.},
``BIG-SPARC,''
(2024); arXiv:2411.13329.

\bibitem{TrivediScherrer2025}
R.~Trivedi and R.~J.~Scherrer,
``Dark Matter from Holography,''
(2025); arXiv:2511.10617.

\end{thebibliography}

\end{document}
